\documentclass[11pt,a4paper]{article}
\usepackage[italian]{babel}
\usepackage[margin=2.5cm]{geometry}
\usepackage{parskip}
\usepackage{amsmath, amssymb, amsthm}
\usepackage{booktabs}
\usepackage{array}
\usepackage{xcolor}
\usepackage{needspace}
\usepackage{enumitem}
\usepackage{listings}
\usepackage{tikz}
\usepackage[most]{tcolorbox}
\usepackage[hidelinks]{hyperref}
\usetikzlibrary{decorations.pathreplacing, shapes.geometric, arrows.meta, positioning}

% Niente vedove/orfane: i paragrafi non lasciano righe isolate a fine/inizio pagina
\widowpenalty=10000
\clubpenalty=10000

% Elenchi più compatti
\setlist{itemsep=2pt, parsep=0pt, topsep=4pt, partopsep=0pt}

% --- Riquadri con bordo nero, senza numero (si spezzano tra due pagine) ---
% Uso: \begin{definizione} ... \end{definizione}
%      \begin{teorema}[Nome] ... \end{teorema}  (il nome è facoltativo)
\tcbset{nota/.style={enhanced jigsaw, breakable, colback=white,
  colframe=black, boxrule=0.5pt, sharp corners}}
\NewTColorBox{definizione}{}{nota,
  before upper={\textbf{Definizione.}\ }}
\NewTColorBox{teorema}{o}{nota,
  before upper={\textbf{Teorema\IfValueT{#1}{ (#1)}.}\ }}
\NewTColorBox{proposizione}{o}{nota,
  before upper={\textbf{Proposizione\IfValueT{#1}{ (#1)}.}\ }}
\NewTColorBox{proprieta}{o}{nota,
  before upper={\textbf{Proprietà\IfValueT{#1}{ (#1)}.}\ }}
\NewTColorBox{assioma}{o}{nota, boxrule=1pt,
  before upper={\textbf{Assioma\IfValueT{#1}{ (#1)}.}\ }}

% --- Ambienti semplici (senza riquadro) ---
% Il titolo sta in cima al blocco, anche se il contenuto inizia con un riquadro o
% con codice; e non resta mai da solo in fondo a una pagina.
% Uso: \begin{esempio}[Titolo facoltativo] ... \end{esempio}
\newcommand{\newrunin}[2]{%
  \NewDocumentEnvironment{#1}{o}{%
    \par\Needspace{4\baselineskip}\addvspace{\medskipamount}\noindent
    \textbf{#2}\IfValueT{##1}{ (##1)}\textbf{.}\ \ignorespaces}%
  {\par\addvspace{\medskipamount}}}
\newrunin{esempio}{Esempio}
\newrunin{esercizio}{Esercizio}
\newrunin{osservazione}{Osservazione}

% V e F nelle tavole di verità
\newcommand{\V}{\textsf{V}}
\newcommand{\F}{\textsf{F}}

% --- Codice C ---
% Uso: \begin{codice} ... \end{codice}      codice C numerato
%      \begin{comando} ... \end{comando}    comandi da terminale
%      \begin{risultato} ... \end{risultato} output di un programma
%      \lstinline|printf("ciao");|           codice dentro al testo
\lstdefinestyle{c}{
  language=C,
  basicstyle=\ttfamily\small,
  keywordstyle=\color{blue}\bfseries,
  commentstyle=\color{gray}\itshape,
  stringstyle=\color{red!70!black},
  numbers=left,
  numberstyle=\tiny\color{gray},
  numbersep=8pt,
  columns=flexible,
  keepspaces=true,
  breaklines=true,
  showstringspaces=false,
  tabsize=4,
  morekeywords={include, define},
  % lettere accentate nei commenti
  literate={à}{{\`a}}1 {è}{{\`e}}1 {é}{{\'e}}1 {ì}{{\`i}}1
           {ò}{{\`o}}1 {ù}{{\`u}}1 {È}{{\`E}}1 {À}{{\`A}}1
}
\lstset{style=c}
\newtcblisting{codice}{listing only, nota, left=6mm,
  listing options={style=c}}
\newtcblisting{comando}{listing only, nota,
  listing options={basicstyle=\ttfamily\small, numbers=none, language={},
    columns=flexible, keepspaces=true, breaklines=true}}
\newtcblisting{risultato}{listing only, enhanced jigsaw, breakable,
  colback=black!5, colframe=black, boxrule=0.5pt, sharp corners,
  listing options={basicstyle=\ttfamily\small, numbers=none, language={},
    columns=flexible, keepspaces=true, breaklines=true}}

% --- Diagrammi di flusso (simboli standard) ---
\tikzset{
  terminale/.style = {ellipse, draw, minimum width=2.5cm, minimum height=0.9cm},
  io/.style        = {trapezium, trapezium left angle=70, trapezium right angle=110,
                      draw, minimum width=2.5cm, minimum height=0.9cm},
  processo/.style  = {rectangle, draw, minimum width=2.5cm, minimum height=0.9cm},
  decisione/.style = {diamond, draw, aspect=2, minimum width=2.5cm},
  freccia/.style   = {thick, -{Stealth}}
}

\title{Teoria degli insiemi}
\author{Marco Cerbella}
\date{Lezione 1 -- 24 settembre 2026}

\begin{document}
\maketitle

\section{Insiemi e loro costruzione}

\begin{definizione}
Un \emph{insieme} è una collezione di oggetti, detti \emph{elementi},
accomunati da qualche proprietà.
\end{definizione}

Un insieme può essere definito nei seguenti modi:
\begin{enumerate}
    \item \textbf{per caratteristica}: $\{n \in \mathbb{N} \mid n \leq 4\}$;
    \item \textbf{per elencazione}: $\{0, 1, 2, 3, 4\}$;
    \item \textbf{con i diagrammi di Eulero-Venn}.
\end{enumerate}

Tutti gli insiemi che si considerano sono sottoinsiemi di un insieme più
vasto, detto \emph{universo} e indicato con $U$. Due insiemi appartenenti
allo stesso universo si possono sempre confrontare tra loro.

\begin{definizione}
Dati due insiemi $A$ e $B$, si dice che $A$ è un \emph{sottoinsieme} di $B$,
e si scrive $A \subseteq B$, se ogni elemento di $A$ è anche elemento di $B$.
\end{definizione}

Se esiste $b \in B$ con $b \notin A$ (cioè $A \subseteq B$ ma $A \neq B$) si dice che $A$ è un \emph{sottoinsieme proprio} di $B$ e si scrive $A \subsetneq B$.

Due insiemi sono uguali, $A = B$, se $A \subseteq B$ e $B \subseteq A$,
cioè se hanno esattamente gli stessi elementi.

Tra i sottoinsiemi di un insieme $A$ si considera anche l'\emph{insieme vuoto},
cioè l'insieme privo di elementi, indicato con il simbolo $\emptyset$.

\section{Intersezione}

\begin{definizione}
L'\emph{intersezione} di $A$ e $B$ è l'insieme degli elementi che
appartengono sia ad $A$ sia a $B$:
\[
A \cap B = \{x \mid x \in A \text{ e } x \in B\}.
\]
\end{definizione}

\begin{esempio}
Se $A = \{a, b, c, d\}$ e $B = \{b, d, f\}$, allora $A \cap B = \{b, d\}$.
\end{esempio}

L'intersezione gode della proprietà commutativa: $A \cap B = B \cap A$.
Altre proprietà:
\begin{enumerate}
    \item se $A = B$ allora $A \cap B = A = B$;
    \item se $A$ e $B$ non hanno elementi in comune allora $A \cap B = \emptyset$ ($A$ e $B$ si dicono \emph{disgiunti});
    \item se $A \subseteq B$ allora $A \cap B = A$;
    \item $A \cap \emptyset = \emptyset$ e $A \cap U = A$.
\end{enumerate}
\section{Unione}

\begin{definizione}
L'\emph{unione} di $A$ e $B$ è l'insieme degli elementi che appartengono
ad $A$, a $B$ o a entrambi:
\[
A \cup B = \{x \mid x \in A \text{ o } x \in B\}.
\]
\end{definizione}

\begin{esempio}
Se $A = \{a, b, c, d\}$ e $B = \{b, d, f\}$, allora
$A \cup B = \{a, b, c, d, f\}$.
\end{esempio}

Anche l'unione gode della proprietà commutativa: $A \cup B = B \cup A$.
Altre proprietà:
\begin{enumerate}
    \item se $A = B$ allora $A \cup B = A = B$;
    \item se $A \neq \emptyset$ oppure $B \neq \emptyset$ allora $A \cup B \neq \emptyset$;
    \item se $A \subseteq B$ allora $A \cup B = B$;
    \item $A \cup \emptyset = A$ e $A \cup U = U$.
\end{enumerate}

\begin{esempio}
Se $A \subseteq B$ allora $A \cup B = B$: tutti gli elementi di $A$ sono già in $B$, quindi unirli non aggiunge nulla.
\end{esempio}

\section{Differenza tra insiemi}

\begin{definizione}
La \emph{differenza} tra $A$ e $B$ è l'insieme degli elementi che
stanno in $A$ ma non in $B$:
\[
A \setminus B = \{x \mid x \in A \text{ e } x \notin B\}.
\]
\end{definizione}

\begin{esempio}
Se $A = \{a, b, c, d\}$ e $B = \{b, d, f\}$, allora
$A \setminus B = \{a, c\}$ e $B \setminus A = \{f\}$.
\end{esempio}

La differenza \textbf{non} gode della proprietà commutativa: in generale
$A \setminus B \neq B \setminus A$ (lo mostra l'esempio precedente).

Note utili:
\begin{enumerate}
    \item $A \setminus A = \emptyset$;
    \item $A \setminus \emptyset = A$;
    \item se $A \subseteq B$, allora $A \setminus B = \emptyset$
          (tutti gli elementi di $A$ stanno anche in $B$, quindi non ne
          resta nessuno);
    \item $A \setminus B = A \cap B^c$ (togliere $B$ da $A$ equivale a
          intersecare $A$ con il complementare di $B$).
\end{enumerate}

\section{Complementare}

\begin{definizione}
Il \emph{complementare} di $A$ è l'insieme di tutti gli elementi di $U$
che non appartengono ad $A$:
\[
A^c = \{x \in U \mid x \notin A\} = U \setminus A.
\]
\end{definizione}

\begin{esempio}
Se $U = \{a, b, c, d, e, f\}$, $A = \{a, b, c, d\}$ e $B = \{b, d, f\}$,
allora $A^c = \{e, f\}$ e $B^c = \{a, c, e\}$.
\end{esempio}

Il \emph{complementare di $A$ rispetto a $B$}, con $A \subseteq B$, si
ottiene usando $B$ al posto dell'universo $U$:
\[
\complement_B A = B \setminus A.
\]
Va sempre indicato esplicitamente: $A^c$ senza altre indicazioni è sempre
il complementare rispetto a $U$.

\begin{esempio}
Se $U = \{1, 2, \dots, 10\}$, $B = \{1, 2, 3, 4, 5\}$ e $A = \{2, 4\}$
(quindi $A \subseteq B$), allora
\[
\complement_B A = \{1, 3, 5\}, \qquad A^c = \{1, 3, 5, 6, 7, 8, 9, 10\}.
\]
\end{esempio}

Proprietà: $(A^c)^c = A$, \quad $A \cup A^c = U$, \quad
$A \cap A^c = \emptyset$, \quad $U^c = \emptyset$, \quad $\emptyset^c = U$.

\begin{teorema}[Leggi di De Morgan]
\begin{enumerate}
    \item $(A \cup B)^c = A^c \cap B^c$: il complementare dell'unione
          è l'intersezione dei complementari;
    \item $(A \cap B)^c = A^c \cup B^c$: il complementare
          dell'intersezione è l'unione dei complementari.
\end{enumerate}
\end{teorema}

Regola pratica: il complementare ``entra'' nella parentesi e unione e
intersezione si scambiano.

\begin{esempio}
Con $U = \{a, b, c, d, e, f\}$, $A = \{a, b, c, d\}$, $B = \{b, d, f\}$:
$(A \cup B)^c = \{a,b,c,d,f\}^c = \{e\}$ e
$A^c \cap B^c = \{e, f\} \cap \{a, c, e\} = \{e\}$.
\end{esempio}

\section{Prodotto cartesiano}

\begin{definizione}
Una \emph{coppia ordinata} $(a, b)$ è formata da due componenti prese in un
ordine preciso. Due coppie sono uguali se hanno le stesse componenti nello
stesso ordine:
\[
(a, b) = (c, d) \iff a = c \text{ e } b = d.
\]
Ad esempio $(1, 2) \neq (2, 1)$, mentre $\{1, 2\} = \{2, 1\}$.
\end{definizione}

\begin{definizione}
Il \emph{prodotto cartesiano} di $A$ e $B$ è l'insieme di tutte e sole le
coppie ordinate con la prima componente in $A$ e la seconda componente
in $B$:
\[
A \times B = \{(a, b) \mid a \in A \text{ e } b \in B\}.
\]
\end{definizione}

\begin{esempio}
Se $A = \{1, 2, 3\}$ e $B = \{4, 5, 6\}$, allora
\[
A \times B = \{(1,4), (1,5), (1,6), (2,4), (2,5), (2,6), (3,4), (3,5), (3,6)\}.
\]
\end{esempio}

Il prodotto cartesiano \textbf{non} gode della proprietà commutativa:
in generale $A \times B \neq B \times A$.

Note utili:
\begin{enumerate}
    \item se $A$ ha $m$ elementi e $B$ ha $n$ elementi, allora
          $A \times B$ ha $m \cdot n$ elementi (nell'esempio $3 \cdot 3 = 9$);
    \item $A \times \emptyset = \emptyset$;
    \item $A \times A$ si scrive direttamente $A^2$.
\end{enumerate}

\textbf{Il piano cartesiano}: $\mathbb{R}^2 = \mathbb{R} \times \mathbb{R}$
è l'insieme di tutte le coppie $(x, y)$ di numeri reali. Ogni coppia
individua un punto del piano ($x$ ascissa, $y$ ordinata).

\end{document}
